% ============================================================
% §5.1.1 Tempo Drift — quantitative results content
% Replace each placeholder block in the draft with the
% corresponding paragraph below.
% ============================================================

\subsubsection{Tempo Drift}
\label{sec:results-tempo}

Based on Research Question~\ref{rq:1.1} and Hypothesis~\ref{hyp:2}.
An 8\% tempo character change was applied relative to each participant's
individual baseline cadence over one minute, with 15-second drift phases.

% ------------------------------------------------------------------
\paragraph{Initial results.}

Figure~\ref{fig:initial-cadence} shows the cadence trajectory as a percentage
change from individual baseline for all 23 sessions.
Across the cohort, the median change at the end of the trial was $-2.42\%$
(IQR $[-7.5,\; +2.6]$\%), with high variability between participants.
Neither the speeding-up group ($n = 14$; median $-0.53\%$;
Wilcoxon $W = 42$, $p = 0.54$, $r = 0.16$)
nor the slowing-down group ($n = 9$; median $-3.32\%$;
Wilcoxon $W = 13$, $p = 0.30$, $r = 0.35$)
showed a statistically significant shift from baseline when assessed
at the trial endpoint.
A binomial test on the rate of correct-direction endpoint change
(12/22 sessions, 55\%) did not differ significantly from chance
($p = 0.83$).

However, a paired comparison of trajectory-based influence scores
between tempo and weight sessions ($n = 22$ paired sessions) reveals
a highly significant difference: tempo sessions produced substantially
higher influence scores (median $0.46$) than weight sessions,
in which no manipulation was applied (median $0.004$;
Wilcoxon $W = 6$, $p < 0.0001$).
This indicates that while group-level endpoint cadence shifts did not reach
significance, the \emph{shape} of cadence trajectories --- whether cadence
tracks the ramp-in, sustains through the plateau, or responds to the
ramp-out --- is significantly more structured in tempo sessions than in
the null condition.
The disparity between endpoint and trajectory analyses reflects the
high within-session cadence variability visible in
Figure~\ref{fig:initial-cadence}; natural drift during free walking
is sufficient to produce apparent correct-direction endpoint changes
in weight sessions at a similar or higher rate than tempo sessions,
making trajectory shape a more discriminating measure than endpoint
alone.

% ------------------------------------------------------------------
\paragraph{Test setup categorisation.}

Figure~\ref{fig:avg-trajectories} shows the mean trajectory for each of the
four test-setup groups.
Participants who completed the tempo condition first showed a higher rate of
correct-direction responses (6/9, 67\%) than those who completed it second
(6/14, 43\%).
The speeding-up, tempo-first group ($n = 3$) produced the largest mean
upward trend (median final $+1.9\%$).
The speeding-up, weight-first group ($n = 11$) moved on average in the
opposite direction ($-1.7\%$).
Slowing-down participants were more directionally consistent
($-0.3\%$ and $-2.2\%$ for tempo-first and weight-first respectively),
though remained well below the 8\% target in both cases.

% ------------------------------------------------------------------
\paragraph{Behaviour categorisation.}

Figure~\ref{fig:behaviour-categories} summarises the distribution across
the four behavioural categories.
Eight of 23 sessions (35\%) showed \textit{Consistent influence},
including both clear and partial directional tracking throughout the trial.
A further eight (35\%) showed \textit{Impact only from drift phases},
responding to the ramp-in or ramp-out transitions but not the sustained
plateau, suggesting sensitivity to rate of change rather than level.
One session (4\%) demonstrated \textit{Subconscious correction}.
Six sessions (26\%) showed no clear directional effect in any phase.

% ------------------------------------------------------------------
\paragraph{Systemised categorisation benchmarked to weight-phase data.}

The same trajectory-based pipeline was applied to the 22 available
weight-condition sessions as a false-positive benchmark.
A McNemar test on paired correct-direction classifications found that
the discordant pairs significantly favoured the weight condition over
tempo ($\chi^2 = 4.17$, $p = 0.041$), confirming that simple
endpoint direction is an unreliable indicator of tempo-driven change
and that natural cadence drift can replicate apparent directional
responses in the absence of any manipulation.

By contrast, the paired Wilcoxon on influence scores --- which
accounts for each participant's natural variability by normalising
against their own weight-session cadence standard deviation ---
demonstrates that tempo sessions produce distinctively structured
trajectories that are not attributable to natural drift
($W = 6$, $p < 0.0001$).
This supports the use of trajectory shape as the primary measure of
manipulation effect, and provides evidence that the tempo manipulation
does influence cadence behaviour at the within-session level, even where
the group-level endpoint shift does not reach significance.
Full false-positive rates per profile category are provided in
Appendix~\ref{app:profiles}.

No significant correlation was found between cadence influence score
and Agency Aggregate ($\rho = -0.13$, $p = 0.56$, $n = 23$).
Agency scores were high overall (median $7.0$, range $3.3$--$9.7$),
suggesting participants perceived a sense of ownership over the audio
regardless of whether their cadence changed measurably.
